\begin{figure}[htbp]
\centering
% Schematic only (no data): a Newton polygon with lower chain S, C1, H, C2 (the adjacency
% S-C1-H-C2 is that of the four essential facets of the one-loop off-shell triangle;
% the 2D shape itself is not a projection of that polytope).
\begin{tikzpicture}[fd,x=0.56cm,y=0.56cm,
  fe poly/.style={draw=gray!60,fill=gray!10,line width=0.6pt,line join=round},
  fe pt/.style={circle,fill=black,draw=none,inner sep=0pt,minimum size=3.4pt},
  fe ipt/.style={circle,fill=gray!55,draw=none,inner sep=0pt,minimum size=3pt},
  fe facet/.style={line width=1.7pt,line cap=round},
  fe normal/.style={draw,line width=0.6pt,-{Stealth[length=4pt,width=3.4pt]}},
  fe sign/.style={circle,fill=white,line width=0.7pt,inner sep=0pt,minimum size=8.5pt,font=\scriptsize\bfseries},
  fe lab/.style={font=\footnotesize}]
  % ======================= (a) Euclidean type =======================
  \begin{scope}
    \draw[lcaxis] (-0.7,-0.7) -- (9.4,-0.7) node[below left=1pt,fe lab] {$x$ 的指数（投影）};
    \draw[lcaxis] (-0.7,-0.7) -- (-0.7,6.1) node[right=1pt,fe lab] {$t$ 的指数};
    \draw[fe poly] (0,3) -- (1,1) -- (3,0) -- (6,0) -- (8,1) -- (8,3) -- (6,5) -- (2,5) -- cycle;
    \draw[fe facet,green!45!black] (0,3) -- (1,1);
    \draw[fe facet,blue!60!black] (1,1) -- (3,0);
    \draw[fe facet,black!65] (3,0) -- (6,0);
    \draw[fe facet,orange!80!black] (6,0) -- (8,1);
    \foreach \p in {(0,3),(1,1),(3,0),(6,0),(8,1),(8,3),(6,5),(2,5),(4,0)} \node[fe pt] at \p {};
    \foreach \p in {(2,2),(3,3),(5,2),(6,3),(4,4),(2,4)} \node[fe ipt] at \p {};
    % inward normals of the lower facets (r_{n+1}>0)
    \draw[fe normal,green!45!black] (0.5,2) -- ++(0.894*1.3,0.447*1.3);
    \draw[fe normal,blue!60!black] (2,0.5) -- ++(0.447*1.3,0.894*1.3);
    \draw[fe normal,black!65] (4.5,0) -- ++(0,1.3) node[above,fe lab,inner sep=1pt] {$r$};
    \draw[fe normal,orange!80!black] (7,0.5) -- ++(-0.447*1.3,0.894*1.3);
    \node[fe lab,text=green!40!black,anchor=east] at (0.38,1.75) {$S$};
    \node[fe lab,text=blue!55!black,anchor=north east,inner sep=1pt] at (1.9,0.42) {$C_1$};
    \node[fe lab,anchor=north,inner sep=1.5pt] at (4.5,-0.02) {$H$};
    \node[fe lab,text=orange!65!black,anchor=north west,inner sep=1pt] at (7.1,0.42) {$C_2$};
    \node at (4.35,-2.75) {(a)\ 欧氏型：系数全正};
  \end{scope}
  % ======================= (b) Minkowski kinematics =======================
  \begin{scope}[xshift=6.0cm]
    \draw[lcaxis] (-0.7,-1.6) -- (9.4,-1.6);
    \draw[lcaxis] (-0.7,-1.6) -- (-0.7,6.1);
    \draw[fe poly] (0,3) -- (1,1) -- (3,0) -- (6,0) -- (8,1) -- (8,3) -- (6,5) -- (2,5) -- cycle;
    \draw[fe facet,green!45!black,line width=1.1pt] (0,3) -- (1,1);
    \draw[fe facet,blue!60!black,line width=1.1pt] (1,1) -- (3,0);
    \draw[fe facet,black!65,line width=1.1pt] (3,0) -- (6,0);
    \draw[fe facet,orange!80!black,line width=1.1pt] (6,0) -- (8,1);
    \foreach \p in {(0,3),(1,1),(8,1),(8,3),(6,5),(2,5)} \node[fe pt] at \p {};
    \foreach \p in {(2,2),(3,3),(5,2),(6,3),(4,4),(2,4)} \node[fe ipt] at \p {};
    % new lower facet after dissection (drawn first, the sign circles cover its ends)
    \draw[red!65!black,line width=1.1pt,dash pattern=on 3.5pt off 2pt] (3,0) -- (4.5,-1) -- (6,0);
    \node[circle,draw=red!65!black,fill=red!15,line width=0.7pt,inner sep=0pt,minimum size=4.4pt] at (4.5,-1) {};
    % opposite-sign monomials on the facet H
    \node[fe sign,draw=blue!60!black] at (3,0) {$+$};
    \node[fe sign,draw=red!65!black] at (6,0) {$-$};
    \draw[red!65!black,line width=0.6pt,{Stealth[length=3.6pt]}-{Stealth[length=3.6pt]}] (3.5,0.6) -- node[above,fe lab,inner sep=1.5pt] {相消} (5.5,0.6);
    \node[fe lab,text=red!60!black,anchor=west] at (5.3,-1.05) {剖分后的新面：隐藏区域（如 Glauber）};
    % inset: parameter space with an interior pinch
    \begin{scope}[shift={(10.3,2.1)}]
      \fill[gray!8] (0,0) rectangle (3.5,3.5);
      \draw[lcaxis] (0,0) -- (3.95,0) node[below,fe lab] {$x_i$};
      \draw[lcaxis] (0,0) -- (0,3.95) node[left,fe lab] {$x_j$};
      \draw[red!65!black,line width=0.8pt] plot[domain=0.55:3.5,samples=40,smooth] (\x,{1.9/\x});
      \draw[red!65!black,line width=1pt] (1.6,1.1875) ++(-0.24,-0.24) -- ++(0.48,0.48) ++(0,-0.48) -- ++(-0.48,0.48);
      \node[fe lab,text=red!60!black,anchor=south west,inner sep=1pt] at (1.85,1.3) {钳夹};
      \node[fe lab,align=center,anchor=north] at (1.75,-0.5) {参数空间：内部钳夹，\\而非端点奇异};
    \end{scope}
    \node at (6.6,-2.75) {(b)\ Minkowski 运动学：异号项（如三圈非平面 $2\to2$）};
  \end{scope}
\end{tikzpicture}
\caption{区域方法的几何（示意）：(a)~当 $\mathcal U+\mathcal F$ 的系数全为正（欧氏型）时，区域对应 Newton 多面体的本质下表面（$S$、$C_1$、$H$、$C_2$，颜色同约化图中的软、共线与硬子图），法向量 $r$（$r_{n+1}>0$）给出各 Feynman 参数的标度 $x_i\to t^{r_i}x_i$；
(b)~Minkowski 运动学中异号项的相消在参数空间内部产生钳夹，标准的下表面算法漏掉相应区域，沿奇异轨迹剖分后出现新的面（红色虚线）——“隐藏区域”\cite{Gardi:2024axt}。
“区域 $=$ 下表面”作为一般陈述在 Minkowski 运动学中不成立。}
\label{fig:front-regions-polytope}
\end{figure}
